\chapter{Descrição da Proposta}
\label{sec:metodologia}

Este capítulo apresenta a proposta desta pesquisa, detalhando o problema investigado, a abordagem adotada, o desenho experimental e os elementos que compõem a solução proposta.

\section{Visão Geral da Proposta}
\label{subsec:visao_proposta}

\textbf{Ideias a abordar:}
\begin{itemize}
    \item Reapresentar brevemente o problema:
    \begin{itemize}
        \item uso crescente de LLMs no estudo para PMP
        \item incerteza sobre confiabilidade
    \end{itemize}
    \item Ideia central da proposta:
    \begin{itemize}
        \item avaliar sistematicamente o desempenho de LLMs
    \end{itemize}
    \item Estratégia geral:
    \begin{itemize}
        \item uso de questões PMP
        \item aplicação de diferentes prompts
        \item análise quantitativa e qualitativa
    \end{itemize}
    \item Resultado esperado:
    \begin{itemize}
        \item compreensão dos limites e capacidades dos LLMs
    \end{itemize}
\end{itemize}




\section{Questões de Pesquisa}
\label{subsec:qp_proposta}

\textbf{Ideias a abordar:}
\begin{itemize}
    \item Apresentar claramente as QPs:
    \begin{itemize}
        \item QP1: Acurácia dos LLMs
        \item QP2: Variação por domínio/tipo de questão
        \item QP3: Padrões de erro
        \item QP4 (opcional): Impacto do prompting
    \end{itemize}
    \item Explicar o papel de cada QP no estudo
\end{itemize}




\section{Modelo Conceitual da Proposta}
\label{subsec:modelo_conceitual}

\textbf{Ideias a abordar:}
\begin{itemize}
    \item Definir os principais elementos do estudo:
    \begin{itemize}
        \item Entrada:
        \begin{itemize}
            \item Questões PMP
            \item Estratégias de prompt
        \end{itemize}
        \item Processo:
        \begin{itemize}
            \item Execução nos LLMs
        \end{itemize}
        \item Saída:
        \begin{itemize}
            \item Respostas dos modelos
            \item Justificativas
        \end{itemize}
    \end{itemize}
    \item Variáveis:
    \begin{itemize}
        \item Independentes:
        \begin{itemize}
            \item Prompt
            \item Tipo de questão
            \item Domínio PMP
        \end{itemize}
        \item Dependentes:
        \begin{itemize}
            \item Acurácia
            \item Consistência
            \item Tipo de erro
        \end{itemize}
    \end{itemize}
    \item Pode incluir um diagrama (fluxo ou arquitetura)
\end{itemize}



\section{Desenho Experimental}
\label{subsec:desenho_exp}

\textbf{Ideias a abordar:}
\begin{itemize}
    \item Estratégia experimental:
    \begin{itemize}
        \item execução controlada de questões
        \item isolamento entre execuções
    \end{itemize}
    \item Configuração:
    \begin{itemize}
        \item cada questão executada separadamente
        \item ausência de contexto acumulado
    \end{itemize}
    \item Repetições:
    \begin{itemize}
        \item execução múltipla para avaliar consistência
    \end{itemize}
    \item Comparações:
    \begin{itemize}
        \item entre prompts
        \item entre tipos de questão
        \item entre domínios PMP
    \end{itemize}
\end{itemize}




\section{Estratégias de Prompting}
\label{subsec:prompting_proposta}

\textbf{Ideias a abordar:}
\begin{itemize}
    \item Definir os tipos de prompts utilizados:
    \begin{itemize}
        \item Zero-shot
        \item Prompt com explicação (CoT simplificado)
        \item Prompt com grounding (opcional)
    \end{itemize}
    \item Justificar escolha:
    \begin{itemize}
        \item avaliar impacto no desempenho
    \end{itemize}
    \item Definir estrutura das respostas esperadas:
    \begin{itemize}
        \item resposta final
        \item justificativa (quando aplicável)
    \end{itemize}
\end{itemize}




\section{Dataset e Caracterização das Questões}
\label{subsec:dataset_proposta}

\textbf{Ideias a abordar:}
\begin{itemize}
    \item Origem das questões:
    \begin{itemize}
        \item simulados, bancos de questões, ou construção própria
    \end{itemize}
    \item Quantidade de questões
    \item Classificação das questões:
    \begin{itemize}
        \item por domínio (People, Process, Business)
        \item por tipo (objetiva, múltipla, situacional)
    \end{itemize}
    \item Garantia de qualidade:
    \begin{itemize}
        \item validação do gabarito
    \end{itemize}
\end{itemize}



\section{Métricas de Avaliação}
\label{subsec:metricas}

\textbf{Ideias a abordar:}
\begin{itemize}
    \item Métrica principal:
    \begin{itemize}
        \item acurácia
    \end{itemize}
    \item Métricas complementares:
    \begin{itemize}
        \item desempenho por domínio
        \item desempenho por tipo de questão
        \item consistência entre execuções
    \end{itemize}
    \item Avaliação qualitativa:
    \begin{itemize}
        \item análise de justificativas
        \item categorização de erros
    \end{itemize}
\end{itemize}


\section{Análise dos Resultados}
\label{subsec:analise}

\textbf{Ideias a abordar:}
\begin{itemize}
    \item Análise quantitativa:
    \begin{itemize}
        \item comparação entre cenários
        \item estatísticas descritivas
    \end{itemize}
    \item Análise qualitativa:
    \begin{itemize}
        \item identificação de padrões de erro
        \item interpretação dos comportamentos dos modelos
    \end{itemize}
    \item Relacionamento com as QPs
\end{itemize}



\section{Considerações Finais do Capítulo}
\label{subsec:proposta_final}

\textbf{Ideias a abordar:}
\begin{itemize}
    \item Síntese da abordagem proposta
    \item Reforço do diferencial:
    \begin{itemize}
        \item análise integrada (prompt + domínio + tipo)
    \end{itemize}
    \item Viabilidade do estudo
    \item Transição para próximos capítulos:
    \begin{itemize}
        \item implementação e resultados
    \end{itemize}
\end{itemize}